% ======================================================================
% FIGURA: Arquitetura RAG Odontológico (Capítulo 23)
% ======================================================================
\begin{figure}[H]
\centering
\begin{tikzpicture}[
    node distance=1.5cm and 2cm,
    box/.style={rectangle, draw=blueaccent, fill=blueaccent!15, minimum width=2.5cm, minimum height=0.8cm, font=\footnotesize, rounded corners=3pt, align=center},
    db/.style={cylinder, draw=cyanaccent, fill=cyanaccent!15, minimum width=2cm, minimum height=1.2cm, font=\footnotesize, shape border rotate=90, align=center},
    arrow/.style={-latex, thick},
]
% Usuário
\node[box, fill=codegreen!20] (user) {Usuário\\Pergunta clínica};

% Retrieval
\node[box, right=of user] (embed) {Embedding\\Transformer\\768-d};
\node[db, right=of embed] (vectordb) {Vector DB\\100 artigos\\ChromaDB};

% Generation
\node[box, below=of embed] (retrieve) {Top-K\\Recuperação\\k=5};
\node[box, right=of retrieve] (llm) {LLM\\GPT-4/LLaMA\\Geração};
\node[box, right=of llm, fill=codegreen!20] (answer) {Resposta\\com DOI\\+ Confiança};

% Setas
\draw[arrow] (user) -- (embed) node[midway, above, font=\tiny] {Query};
\draw[arrow] (embed) -- (vectordb) node[midway, above, font=\tiny] {similaridade};
\draw[arrow] (vectordb) -- (retrieve) node[midway, right, font=\tiny] {trechos};
\draw[arrow] (retrieve) -- (llm) node[midway, above, font=\tiny] {contexto};
\draw[arrow] (llm) -- (answer) node[midway, above, font=\tiny] {resposta};
\draw[arrow, bend left=20] (answer.south) to (user.south);

% Título
\node[font=\small\bfseries, above=0.5cm of embed] {RAG Odontológico: Recuperação + Geração};

\end{tikzpicture}
\caption{Arquitetura de Retrieval-Augmented Generation (RAG) para consulta odontológica. A pergunta do usuário é convertida em embedding (768-d), comparada com a base vetorial de 100 artigos científicos, e os 5 trechos mais relevantes são usados como contexto pelo LLM para gerar uma resposta com citação DOI.}
\label{fig:rag-architecture}
\end{figure}
